% !TEX root = ../main.tex
\chapter*{Wstęp}
\addcontentsline{toc}{chapter}{Wstęp}

Rekonstrukcja trójwymiarowych scen ze zdjęć jest jednym z~ważniejszych
problemów współczesnej grafiki komputerowej. Metody oparte na geometrii
wielowidokowej, takie jak Structure from Motion i~Multi-View Stereo,
dostarczały jedynie przybliżonych reprezentacji pozbawionych
fotorealistycznego wyglądu. Krokiem naprzód w~tej dziedzinie była nowa
technika neuronowych pól luminancji (NeRF), jednak jej praktyczne
zastosowanie ograniczały wielogodzinny czas treningu oraz brak możliwości
renderowania w~czasie rzeczywistym, co sprawiło, że trzeba było szukać innej
metody.

Opublikowana w~2023 roku metoda 3D Gaussian Splatting (3DGS) rozwiązuje
problemy z~treningiem i~real-time renderingiem. Reprezentuje ona scenę jako
zbiór trójwymiarowych Gaussianów i~renderuje ją za pomocą różniczkowalnego
rasteryzatora kafelkowego na GPU. Metoda ta osiąga jakość wizualną
porównywalną z~NeRF przy kilkudziesięciu klatkach na sekundę. Krótko po
publikacji 3DGS znalazł zastosowanie w~produkcji filmowej, wizualizacji
architektonicznej, a~także w~badaniach nad robotyką i~autonomicznymi
pojazdami.

Dynamiczny rozwój metody ujawnił problemy w~dostępnych narzędziach do jej
wizualizacji. Istniejące rozwiązania zmuszają użytkownika do wyboru między
narzędziami analitycznymi, narzędziami do edycji jak i~UX. Narzędzia
analityczne, jak splatviz, wymagają karty graficznej NVIDIA z~obsługą CUDA
i~oferują nieprzyjazny użytkownikowi interfejs. Narzędzia produkcyjne, jak
SuperSplat oparty na WebGL, cechują się dojrzałym UX i~możliwością edycji
scen, ale są pozbawione jakichkolwiek narzędzi do analizy scen.

Implementacje niezależne od CUDA, jak 3DGS.cpp oparte na Vulkanie, pozostają
na poziomie bazowym, bez możliwości debugowania czy przeglądania statystyk
sceny, jak i~bez możliwości edycji. Żadne z~dostępnych narzędzi nie łączy
niezależności od CUDA z~podstawową analityką, edycją i~przyjaznym UI.

Celem niniejszej pracy jest zaprojektowanie i~implementacja systemu
wizualizacji scen trójwymiarowych opartego na metodzie 3DGS, działającego
niezależnie od technologii CUDA, z~wykorzystaniem języka C++, graficznego
API OpenGL i~wprowadzonych w~wersji 4.3 OpenGL-a compute shaderów.

Praca ta obejmuje analizę istniejących rozwiązań, implementację całego
potoku graficznego -- od wczytywania splatów do rasteryzacji Gaussianów --
oraz testy wydajnościowe.
